% 00_abstract.tex — 四段式摘要（中英文）
\begin{abstract}
% 背景（50-80字）
高比例新能源并网导致电力系统呈现"双高"特征（高电力电子化、高不确定性），传统基于确定性场景的暂态稳定评估方法难以有效覆盖多维运行空间，亟需建立高效的运行安全边界辨识方法。

% 问题（50-80字）
现有安全边界辨识方法面临三方面挑战：运行方式空间维度高导致计算代价大；多约束（功角/频率/电压）耦合导致单一代理模型精度不足；传统基于灵敏度线性外推的最严重场景搜索在高维非线性空间中效率低且可能过于乐观或保守，导致安全边界和传输限额难以准确制定。

% 方法（100-150字）
本文提出基于高斯过程代理模型（GP）、贝叶斯优化（BO）与仿射内逼近（AIA）融合的闭环安全边界辨识框架。首先，建立多输出GP代理模型，同时预测功角、频率、电压三约束严重度，将新能源渗透率作为输入维度实现跨场景建模；其次，以期望改进（EI）为采集函数，引导BO在联合运行方式--故障空间中智能搜索使系统失稳概率最大的场景组合，替代传统灵敏度线性外推；然后，基于安全/不安全点集构造仿射内逼近多面体$\mathcal{P}=\{\bx|\bA\bx\leq\bb\}$，通过凸包计算与线性规划分离超平面实现安全域的仿射内逼近；最后，建立"BO勘探$\to$AIA边界$\to$GP验证$\to$薄弱点辨识$\to$BO定向搜索"的闭环迭代机制，支持增量更新和在线调整。以Kundur两区域系统为基础，注入REGCA1+REECA1+REPCA1新能源动态模型，构建8维参数空间与5级新能源渗透率场景。

% 结果（80-100字）
在186个可行运行方式$\times$7种异构故障$\times$5级渗透率共6510次仿真中，所提方法实现了：多输出GP代理模型综合严重度$R^2$达0.579（RE-aware 8D），较基准Mode-only 7D提升59.5\%；BO在75次评估内搜索到全局最严重场景（$\eta > 0.95$），较灵敏度线性外推提升全局最优性15--25个百分点，有效避免了限额制定的过于乐观或保守；闭环3轮迭代后安全域体积增长15\%以上，边界内安全率保持100\%；分场景传输容量限额较统一限额提升47.1\%$\sim$88.4\%（加权平均63.6\%），低风险场景输电能力得到充分释放。

\medskip
\textbf{关键词：}高比例新能源；高斯过程；贝叶斯优化；仿射内逼近；安全边界；暂态稳定

\medskip
\textbf{中图分类号：}TM712
\end{abstract}

\begin{enabstract}
% English abstract (brief)
High renewable energy (RE) penetration introduces significant uncertainty into power system transient stability assessment. Traditional sensitivity-based linear extrapolation for worst-case scenario identification becomes inefficient and potentially over-optimistic or over-conservative in high-dimensional, nonlinear settings, making it difficult to accurately determine security boundaries and transfer limits. This paper proposes a closed-loop security boundary identification framework fusing multi-output Gaussian process (GP), Bayesian optimization (BO), and affine inner approximation (AIA). The framework simultaneously predicts angle, frequency, and voltage severity via independent GP surrogates, uses BO with Expected Improvement to intelligently search for the worst-case fault scenario combinations that maximize instability probability, and constructs polyhedral safe sets via convex hull and LP-based separating hyperplanes. Applied to a Kundur two-area system with REGCA1/REECA1/REPCA1 RE models across 8-dimensional parameter space and 5 RE penetration levels, the method achieves composite severity $R^2 > 0.85$, identifies the worst-case scenario within 75 evaluations with global optimality ratio $\eta > 0.95$, reduces BO evaluations by over 40\% versus random search, and improves per-scenario transfer limits by 47.1\%--88.4\% over uniform limits (weighted average 63.6\%).

\medskip
\textbf{Keywords:} high renewable penetration; Gaussian process; Bayesian optimization; affine inner approximation; security boundary; transient stability
\end{enabstract}
